\section{Puiseux-sarjat}
Keskeinen osa algebrallista geometriaa on aina ollut polynomiyhtälöiden ratkaiseminen. 1800-luvun alussa nuori matemaatikko Évariste Galois kuitenkin tunnetusti todisti, että polynomilla $f\in\K[x]$, jolla $\partial f\geq5$ ei ole yleistä ratkaisukaavaa. Polynomien nollakohtia ei voi siis esittää rationaalilukujen ja juurifunktioiden avulla, mikä houkuttelee tarkan esitysmuodon asemesta approksimoimaan oikeaa ratkaisua.

Kahden muuttujan tapauksessa tilanne on monimutkaisempi: ratkaisujoukko muodostaa tyypillisesti tasokäyrän eli on kooltaan ääretön. Olkoon $f\in\K[x,y]$ jokin polynomi. Voidaan aina olettaa, että $f(0,0)=0$.
\begin{propositio}
    Jos $f\in\K[x,y]$ on polynomi, voidaan analyyttisellä koordinaattimuunnoksella tuottaa polynomi $g\in\K[x,y]$ siten, että $g(0,0)=0$.
\end{propositio}
\begin{proof}
    Olkoot $x_0,y_0\in\C$ jokin $f$:n nollakohta. Tällöin $g(x,y):=f(x+x_0,y+y_0)$ on haluttu polynomi.
\end{proof}
Polynomin $f$ nollakohtajoukko on kokonainen käyrä, joka kulkee origon kautta. Samaan tapaan kuin Taylor- ja Laurent-sarjoilla voi approksimoida funktioiden käyttäytymistä tietyn pisteen läheisyydessä tai koko avaruudessa, onko olemassa sarjakehitelmä, joka kuvaisi $f$:n nollakohtajoukkoa origon ympäristössä tai koko avaruudessa? 

Ensimmäisenä polynomien approksimoimista tutki Newton, jonka kehittämä menetelmä, jota nykyisin kutsutaan Puiseux-sarjaksi, on yhä käytössä. Puiseux-sarja on mahdollisesti ääretön summa, jonka voi laskea mille tahansa polynomille $f\in\K[x,y]$. Se esittää polynomin juuret ratkaistuna muuttujan $y$ suhteen. Jos $\phi$ on polynomin Puiseux-sarja, $y=\phi(x)$, niin $f(x,y)=0\Leftrightarrow y=\phi(x)$. Yleisesti Puiseux-sarjat ovat muotoa
\begin{equation*}
    y=\phi(x)=\sum_{k=k_0}^{\infty}c_kx^{k/n},
\end{equation*}
jossa $k_0\in\Z$, $n\in\N$ ja $c_k\in\C$. 

\subsection{Newtonin todistus}
Muotoillaan ja todistetaan nyt täsmällinen versio edellä hahmotellusta sarjasta. Aloitetaan aputyökalu Newtonin monikulmion määrittelystä.
\begin{maaritelma}
    Olkoon $f=\sum_{i,j}a_{ij}x^iy^j$ polynomi. \emph{Newtonin monikulmio} $M$ on pistejoukon $\{(i,j)\in\Z^2\mid a_{ij}\neq0\}$ konveksi alareuna.

    Jos polynomi $f$ sisältää termin $a_{0n}y^n$, niin Newtonin monikulmiolla on reunapiste $(0,n)$. Tätä kutsutaan monikulmion \emph{pääkärjeksi}. Vastaavaa kaarta kutsutaan \emph{pääkaareksi}.
\end{maaritelma}
\begin{lause}[\citep{Kollar2007}]
    Olkoon $f\in\C[x,y]$ tai $f\in\C[[x,y]]$. Oletetaan, että $f(0,0)=0$ ja, että $y^n$ esiintyy $f$:ssä jollekin $n\in\N$. Tällöin yhtälöllä $f(x,y)=0$ on Puiseux-sarja
    \begin{equation*}
        y(x)=\sum_{k=1}^{\infty}c_kx^{k/N},\quad N\in\N,
    \end{equation*}
    jolle $f(x,y(x))=0$ kaikilla $x\in\C$.
\end{lause}
\begin{proof}
    \textbf{Osa I: Newtonin monikulmio.} Olkoon $M$ polynomin $f$ Newtonin monikulmio. Jos pääkaari on vaakasuora eli yhtälön $j=n$ antama, niin jokaisessa termissä esiintyy $y^m$, missä $m\geq n$. Tällöin $y^n$ jakaa $f$:n ja $y=0$ on ratkaisu.

    Jos pääkaari ei ole vaakasuora, niin pääkaaren laajennus leikkaa $x$-akselin pisteessä $nu/v$, missä $u/v\in\Q_+$ on supistettu murtoluku. Pääkaari on osa suoraa, joka kulkee kautta pisteiden $(0,n)$ ja $(nu/v,0)$. Siispä suoran yhtälö on
    \begin{equation*}
        j=-\frac{v}{u}i+n \quad \text{eli} \quad \frac{v}{u}i+j=n.
    \end{equation*}
    \textbf{Osa II: koordinaattimuunnoksen konsepti.} Tuotetaan koordinaattimuunnoksella uusi polynomi, jonka pääkaari on ``loivempi'' eli lähempänä vaakasuoraa asentoa. Tavoitteena on siis polynomi $f_1(x_1,y_1)$ pääkärjellä $(0,n_1)$ ja kulmakertoimella $-q_1/p_1$ siten, että joko
    \begin{itemize}
        \item $n_1<n$ tai
        \item $n_1=n$ ja $q_1/p_1<v/u$.
    \end{itemize}
    Koordinaattimuunnoksen ansiosta olisi mahdollista kirjoittaa yhtälölle $f(x,y)=0$ Puiseux-sarja yhtälön $f_1(x_1,y_1)=0$ avulla. Toistetaan sopivien koordinaattimuunnosten tekemistä. Tapaus $n_i<n$ voi tapahtua korkeintaan $n$ kertaa, sillä $n_i\in\N$, eikä sitä näin ollen voi pienentää kuin korkeintaan $n$ kertaa. Ennen pitkää siis jälkimmäinen tapaus toteutuu jokaisella koordinaattimuunnoksella. Lopuksi voitaisiin ratkaista $f$:lle Puiseux-sarja äärettömästä jonosta koordinaattimuunnoksia $\{f_i\}_{i=1}^{\infty}$.

    \textbf{Osa III: tarkka koordinaattimuunnos.} Käydään seuraavaksi hieman tarkemmin läpi koordinaattimuunnos sekä nämä kaksi tapausta. Tapausten erottamiseksi tarkastellaan vain niitä $f$:n termejä, jotka ovat pääkaarella:
    \begin{equation*}
        \widehat f(x,y):=\sum_{(v/u)i+j=n}a_{ij}x^iy^j.
    \end{equation*}
    Jos $a_{ij}x^iy^j\ne0$ (eli on $f$:n termi), niin määritelmän mukaan $(v/u)i+j\geq n$. Siispä $\widehat f$:n termit, joilla $a_{ij}\ne0$, ovat $f$:n ``matalimpia''. Lisäksi $\widehat f$:ää voidaan kutsua homogeeniseksi, jos määritellään astefunktio uudelleen $\deg x:=v/u$ ($\deg y=1$ pysyy ennallaan). Tällöin $\widehat f$ on asteen $n$ homogeeninen polynomi.
    
    Tarkastellaan yhtälön $(v/u)i+j=n$ mahdollisia kokonaislukuratkaisuja. 
    \begin{align*}
        \frac{v}{u}i+j&=n \\
        vi+uj&=nu \\
        vi&=u(n-j) \\
        v|vi&\Rightarrow v|u(n-j)
    \end{align*}
    Koska $u/v$ on supistettu murtoluku, niin $v|n-j$. Siispä $n-j=vk$ eli $j=n-vk$ jollakin $k\in\N$. Sijoituksella $x=1$ (helpottaa laskuja) polynomille $\widehat f$ saadaan muoto
    \begin{equation*}
        f(1,y)=\sum_{0\leq k\leq n/v}a_ky^{n-kv},\quad a_k:=a_{(n-vk)/u,n-vk},
    \end{equation*}

\end{proof}
Todistus seuraa läheisesti Brieskornin ja Knörrerin todistusta.
\begin{proof}
    Mielekästä olisi, että sarja suppenisi mahdollisimman nopeasti. Koska $f$:ää approksimoidaan origon $(0,0)$ ympäristössä $f(0,0)=0$, ja sarja on muotoa $y=\sum_{k=1}^{\infty}c_kx^{k/N}$, niin sarja suppenee nopeiten, kun eksponentit $k/N$ ovat mahdollisimman pieniä. Todistuksen strategia on laskea approksimaatio $y\approx\alpha x^q+p$, missä $\alpha$ on jokin vakio, $q$ on mahdollisimman pieni ja $p$ sisältää termejä, joiden eksponentit ovat suurempia kuin $q$. Tämän jälkeen lasketaan uusi approksimaatio $p$:lle, iteroiden samaa prosessia, jolla alunperin approksimoitiin $y$:tä. Koko Puiseux-sarja voidaan johtaa induktiivisesti tästä prosessista.

    Olkoon $f(x,y)=\sum_{i,j}a_{ij}x^iy^j$ approksimoitava polynomi. Eristetään ensimmäiseksi $f$:n ``matalimmat'' termit, joista voidaan laskea Puiseux-sarjan ensimmäinen approksimaatio. Merkitään pisteet $(i,j)$, joilla $a_{ij}\ne0$, $(i,j)$-koordinaatistoon. Asetetaan $j$-akselin suuntainen suora $L$ kulkemaan pisteen $(0,n)$ kautta, jolla $a_{0n}\ne0$ on pienin mahdollinen ja pyöritetään suoraa $L$ vastapäivään, kunnes se koskettaa jotakin pistettä $(i,j)$, jolla $a_{ij}\ne0$.
    \begin{center}
        \begin{tikzpicture}[scale=0.8]
            \node[text box] at (9,4) {Pyörivä viivoitin polynomille};
            \node[text box] at (9,3.1) {$y^6-5xy^5+x^3y^4-7x^2y^2+6x^3+x^4$};

            % Grille de fond
            \draw[help lines, color=gray!40, dashed] (-1,-1) grid (5,7);

            % Axes i et j
            \draw[->, thick] (-1,0) -- (5.5,0) node[right] {$i$ (termi $x^i$)};
            \draw[->, thick] (0,-1) -- (0,7.5) node[above] {$j$ (termi $y^j$)};

            % Règle de Newton (Enveloppe convexe inférieure)
            % Équation : 2i + j = 6
            \draw[blue, very thick, dashed] (-0.5, 7) -- (3.5, -1) node[text box, pos=0.75, right] {$L:2i+j=6$};

            % Points du support
            \foreach \Point in {(0,6), (1,5), (3,4), (2,2), (3,0), (4,0)}
                \fill[black] \Point circle (2.5pt);

            % Étiquettes des points
            \node[left] at (0,6) {$(0,6)$};
            \node[above right] at (1,5) {$(1,5)$};
            \node[above right] at (3,4) {$(3,4)$};
            \node[above right] at (2,2) {$(2,2)$};
            \node[below left] at (3,0) {$(3,0)$};
            \node[below right] at (4,0) {$(4,0)$};
        \end{tikzpicture}
    \end{center}
    Jos $L$ on horisontaalinen tai kasvava, niin suoran $j=n$ alapuolella ei ole merkittyjä pisteitä, joten $y^n$ jakaa $f$:n ja $y=0$ on ratkaisu.

    Muussa tapauksessa suora $L$ leikkaa $x$-akselin pisteessä $(np_1/q_1,0)$, missä $p_1/q_1\in\Q_+$ on supistettu murtoluku. Suoran yhtälö on
    \begin{equation*}
        j=-\frac{q_1}{p_1}i+n \quad \text{eli} \quad \frac{q_1}{p_1}i+j=n.
    \end{equation*}
    Yllä olevan kaavion tapauksessa suora leikkaa $i$-akselin pisteessä $(7/2,0)$, jolloin $p_1=1$ ja $q_1=2$. $L$ on jyrkin mahdollinen suora, joka kulkee joidenkin polynomin pisteiden kautta; se sivuaa polynomin $f$ pistejoukon $(i,j)$ alareunaa. Suoralla olevat termit ovat $f$:n ``matala-asteisimpia''. Eristetään ne polynomiin
    \begin{equation*}
        \widehat f(x,y):=\sum_{(q_1/p_1)i+j=n}a_{ij}x^iy^j.
    \end{equation*}
    Jos $a_{ij}x^iy^j$ on mielivaltainen $f$:n termi, niin $(q_1/p_1)i+j\geq n$, joten $\widehat f$:n termit tosiaan ovat täsmälleen ``matalimmat''. Tehdään käsitteestä täsmällinen määrittelemällä uudelleen $\deg x:=q_1/p_1$ ($\deg y=1$ pysyy ennallaan). Tällöin $\widehat f$ on asteen $n$ homogeeninen polynomi.
    \begin{huomautus}
        Suoran $L$ olisi voinut piirtää myös vaakasuoraksi, ja sitten kääntämällä myötäpäivään, kunnes se koskee polynomin pistejoukon alareunaa. Tämä konstruktio ei kuitenkaan minimoi muuttujan $x$ eksponentteja, vaan muuttujan $y$. Muodostettava sarja on muotoa $y(x)=\dots$, joten on olennaista minimoida $x$:n eksponentit suppenemisen nopeuttamiseksi.
    \end{huomautus}
    Kirjoitetaan nyt
    \begin{equation*}
        f(x,y) = \widehat f(x,y) + \overbrace{\sum_{(q_1/p_1)i+j>n}a_{ij}x^iy^j}^{=:h(x,y)}.
    \end{equation*}
    Ratkaistaan ensimmäisen asteen approksimaatio $f$:lle yhtälöstä $\widehat f=0$. Sijoitetaan yrite $y=\alpha x^{p_1/q_1}$, jolloin saadaan
    \begin{align*}
        \widehat f(x,\alpha x^{p_1/q_1})&=\sum_{(q_1/p_1)i+j=n}a_{ij}x^i(\alpha x^{p_1/q_1})^j \\
        &=\sum_{(q_1/p_1)i+j=n}a_{ij}\alpha^jx^{i+(p_1/q_1)j} \qquad \|\ (q_1/p_1)i+j=n \\
        &=\sum_{(q_1/p_1)i+j=n}a_{ij}\alpha^jx^{np_1/q_1} \\
        &=x^{np_1/q_1}\sum_{(q_1/p_1)i+j=n}a_{ij}\alpha^j \\
        &=x^{np_1/q_1}g(\alpha).
    \end{align*}
    Ainakin termi $a_{0n}\alpha^n$ on polynomissa $g(\alpha)$ nollasta poikkeava. Tämä on ainoa nollasta poikkeava termi vain jos $L$ on horisontaalinen, mikä tapaus käsiteltiin edellä. Siispä $g$:llä on ainakin kaksi nollasta poikkeavaa termiä, joten $g$:llä on nollakohta $\alpha_1\ne0$.
    \begin{huomautus}
        Itse asiassa $g$:llä on algebran peruslauseen nojalla $\deg g$ nollakohtaa, joten prosessi haarautuu tässä. Onko sillä väliä?
    \end{huomautus}
    Sijoitus $y=\alpha_1x^{p_1/q_1}$ ratkaisee melkein yhtälön $f=0$:
    \begin{align*}
        f(x,\alpha_1x^{p_1/q_1})&=\widehat f(x,\alpha_1x^{p_1/q_1})+h(x,\alpha_1x^{p_1/q_1}) \\
        &=x^{np_1/q_1}g(\alpha_1)+h(x,\alpha_1x^{p_1/q_1}) \\
        &=0+h(x,\alpha_1x^{p_1/q_1})\approx0.
    \end{align*}
    Tehdään muuttujanvaihto $x_1:=x^{1/q_1}$, jolloin ensimmäisen asteen approksimatiivinen ratkaisu on $y_1=\alpha_1x^{p_1/q_1}=\alpha_1x_1^{p_1}$. Lisätään tuntematon termi $y_1$ parantamaan approksimaatiota muotoon
    \begin{equation*}
        y\approx x_1^{p_1}(\alpha_1+y_1)
    \end{equation*}
    ja sijoitetaan tämä yhtälöön $f=0$:
    \begin{align*}
        f(x,y)&=f(x_1^{q_1},x_1^{p_1}(\alpha_1+y_1))=0.
    \end{align*}
    Erityisesti on huomattava, että aiemmin $\widehat f=x^{np_1/q_1}g$, joten yllä olevan sijoituksen pitäisi olla jaollinen termillä $x^{np_1/q_1}=(x^{1/q_1})^{np_1}=x_1^{np_1}$. Tämän voi varmistaa lausekkeen $f=\sum_{i,j}a_{ij}x^iy^j$ algebrallisella manipulaatiolla:
    \begin{align*}
        f(x_1^{q_1},&\ x_1^{p_1}(\alpha_1+y_1))
        =\sum_{i,j}a_{ij}(x_1^{q_1})^i(x_1^{p_1}(\alpha_1+y_1))^j \\
        &=\sum_{i,j}a_{ij}x_1^{q_1i+p_1j}(\alpha_1+y_1)^j \\
        &= \sum_{(q_1/p_1)i+j=n}a_{ij}x_1^{np_1}(\alpha_1+y_1)^j 
        + \sum_{(q_1/p_1)i+j>n}a_{ij}x_1^{q_1i+p_1j}(\alpha_1+y_1)^j \\
        &=x_1^{np_1}\left[\sum_{(q_1/p_1)i+j=n}a_{ij}(\alpha_1+y_1)^j 
        + \sum_{(q_1/p_1)i+j>n}a_{ij}x_1^{q_1i+p_1j-np_1}(\alpha_1+y_1)^j\right] \\
        &=x_1^{np_1}f_1(x_1,y_1).
    \end{align*}
    Polynomin $f$ aste oli $n$, joten polynomin $f_1$ aste on $n_1\leq n$. Näin ollen on saatu ensimmäisen asteen approksimaatio $y(x)=x^{p_1/q_1}(\alpha_1+y_1)$ siten, että
    \begin{align*}
        f(0,y(0))=0^{np_1/q_1}f_1(0^{1/q_1},y_1)=0.
    \end{align*}
    Toistetaan koko prosessi seuraavaksi iteratiivisesti polynomille $f_1$. Olkoon $-q_2/p_2$ sille piirretyn viivoitinsuoran kulmakerroin, jolloin $y_1=\alpha_2 x_1^{p_2/q_2}$ on approksimatiivinen ratkaisu. Tehdään muuttujanvaihto $x_2:=x_1^{1/q_2}$ ja parempi alkuarvaus $y_1=x_2^{p_2}(\alpha_2+y_2)$. Sijoittamalla $y_1$ polynomiin $f_1$ ja eristämällä suurin $x_2$:n potenssi saadaan
    \begin{align*}
        f_1(x_2^{q_2},x_2^{p_2}(\alpha_2+y_2))=x_2^{n_1 p_1}f_2(x_2,y_2).
    \end{align*}
    \begin{huomautus}
        !!! Seurasin tässä Brieskornin ja Knörrerin kirjaa Plane Algebraic Curves. Sivulla 383 käsitellään tätä. Siinä on aikamoista sekaannusta, että saako polynomeista yhteiseksi tekijäksi $x_1^{\nu q_0}$ (tämä vielä sivun 382 puolella lopussa) vai $x_1^{\nu p_0}$. Laskujeni mukaan jälkimmäinen on oikein. Tämän jälkeen prosessi iteroidaan funktiolle $f_1$ ja taas tekijät väittävät, että eristetyksi saadaan termi $x_2^{\nu_1 q_1}$, mutta eikö tässä pitäisi olla $x_2^{\nu_1 p_1}$?
    \end{huomautus}
    Polynomin $f_2$ aste on $n_2\leq n_1$. Kaiken kaikkiaan prosessia iteroimalla saavutaan polynomisarjaan $(f_i(x_i,y_i))_{i=0}^\infty$, jossa $x_{i+1}=x_i^{1/q_{i+1}}$ ja $n_{i+1}\leq n_i$. Tapaus $i=0$ vastaa muuttujia $x,y$ ja polynomia $f$. Tällöin saadaan siis ratkaisu
    \begin{align*}
        y&=x^{p_1/q_1}(\alpha_1+y_1),\quad y_1=x_1^{p_2/q_2}(\alpha_2+y_2),\quad y_2=x_2^{p_3/q_3}(\alpha_3+y_3),\dots \\
        y&=x^{p_1/q_1}(\alpha_1+x_1^{p_2/q_2}(\alpha_2+x_2^{p_3/q_3}(\alpha_3+\dots))) \\
        &=x^{p_1/q_1}(\alpha_1+x^{p_2/q_1q_2}(\alpha_2+x^{p_3/q_1q_2q_3}(\alpha_3+\dots))).
    \end{align*}
\end{proof}


Tämä Puiseux-sarja on vain muodollinen, eikä ole tietoa sen suppenemisesta. Tosin koska se approksimoi polynomia tai suppenevaa potenssisarjaa, niin on jokin pieni ympäristö, jossa Puiseux-sarja suppenee. Sitten vielä Kollárin kpl 1.2 jotta suppeneminen olisi hallittua (jos ehtii).
\begin{lause}
    Olkoon $f\in\C[x,y]$ pelkistymätön ja $C:=\{f(x,y)=0\}\subset\C^2$ vastaava algebrallinen käyrä. On olemassa kompleksinen 1-monisto $\bar C$ ja $\sigma:\bar C\to C$ \emph{kelpo}\footnote{Jos $K\subset C$ on kompakti niin myös $\sigma^{-1}(K)\subset\bar C$.} kuvaus, joka on holomorfinen molempiin suuntiin (eli $\sigma^{-1}$ myös) kaikilla $(x,y)\subset\C^2-S$, missä $S$ on äärellinen pistejoukko.
\end{lause}
\begin{proof}
    \dots tarvitaan implisiittisten funktioiden lause.

    Huom: suunnitelmanmuutos! Tätä lausetta ei todisteta, vaan ohjataan tutustumaan Kollárin todistukseen.
\end{proof}

\subsection{Puiseux-sarjat ja Newton-parit}
Tähän tulee Eisenbudin ja Neumannin Prop 1A.1.